% page 520 of the facsimile (image p-519 of the scan)
% block 162.6 x 237.7 mm ; left margin 8.0 mm
\pgc{-12.700mm}{0.000mm}{
\l{\cel{3}512.}
\l{}
\l{\sou{saponario}. (\sou{bot.)}\cel{1}Planto \textquotedbl{}kariofilea\textquotedbl{}, sudorifigiva, dfqua}
\l{\cel{3}onu aquo-boliigas la folii por netigar la lanaji, e c.-}
\l{\cel{3}L.\cel{1}\sou{saponaria}. FISL.}
\l{\cel{21}32\cel{2}54\cel{2}18}
\l{\sou{saponino}. Glukosido C\cel{3}H\cel{3}O\cel{2}, kontenata en saponario}
\l{\cel{3}precipue, e di qua la solvuro spumifas quale aquo sapo-}
\l{\cel{3}noza. -}
\l{}
\l{\sou{saponito}. (\sou{mineralogio}).Silikato naturala amorfa, blankatra,}
\l{\cel{3}vicina ye talko e magnezito, ma qua kontenas alumino.}
\l{}
\l{\sou{saporar}. (\sou{netr}.) (\sou{aludante manjajo o drinkajo}).Efikar, per}
\l{\cel{3}sua propraji kemiala, agreable o desagreable a la papili}
\l{\cel{3}di la lango lor mastiko o drinko. - FIS.}
\l{}
\l{\sou{saprofito}. (\sou{mikrobiol}.)Nomo di la mikrobi qui, normale,}
\l{\cel{3}ne povas developar su en organismo vivanta}
\l{\cel{3}e, konseque, nutras su de materii mortinta. -}
\l{}
\l{\sou{sapto}. (\sou{bot.}) Liquido ek la suki quin la planto cherpas ek}
\l{\cel{3}tero per sua radiki, e qua, absorbite da sua parti diversa,}
\l{\cel{3}nutras lu. - DE.}
\l{}
\l{\sou{sarabando}.Olima danso tri-tempa, di qua la ritmo esis grava}
\l{\cel{3}e lenta. - DEFIRS.}
\l{}
\l{\sou{saraceno}. (\sou{bot.}) Genero de poligonacei di qui la flori, grapa,}
\l{\cel{3}ne havas korolo. La farino ek la grani uzesas por facar gale-}
\l{\cel{3}ti, paplo, krespi, e la grani, por nutrar la kavali, la pul-}
\l{\cel{3}tro, e c. - L.\cel{1}\sou{fagopyrus}. - FIS.}
\l{}
\l{\sou{sarcelo.(zool.}) Aquo-ucelo analoga a la anado, ma plu mikra. -}
\l{\cel{3}L.\cel{1}\sou{anas querquedula}. - F.}
\l{}
\l{\sou{sardino}. (\sou{zool.}) Mikra fisho quan onu manjas fresha, o kon-}
\l{\cel{3}servita en aquo saloza, o macerita en oleo. - L.\cel{1}\sou{clupea}}
\l{\cel{3}\sou{sardina}. - DEFIRS.}
\l{}
\l{\sou{sardonika}. Qua donas a la boko karaktero di moko malignala.}
\l{\cel{3}- DEFIRS.}
\l{}
\l{\sou{sargaso}. (\sou{bot.)}\cel{2}Algo tropikala. -L.\cel{1}\sou{sargassum.}\cel{1}-}
\l{}
\l{\sou{sarigo}. (\sou{zool.)}\cel{2}Mamifero ek la kategorio \textquotedbl{}marsupiali\textquotedbl{}, di}
\l{\cel{3}qua la femino havas posho ventrala ube lua yuni finas deve-}
\l{\cel{3}lopar su. - L. didelphys virginiana.- FI.}
\l{}
\l{\sou{sarko}. I. Che la Greki e Romani antiqua : sorto de kofro sen-}
\l{\cel{3}kovrita, en qua la mortinto portesis a la rogo od a la tom-}
\l{\cel{3}bo. - II. Che la moderni : kofro de ligno, de plombo, e c.,}
\l{\cel{3}en qua onu enklozas la mortinti por sepultar li. - DE.}
\l{}
\l{\sou{sarkasmo}. Moko, ironio mordanta. - DEFIRS.}
}
